% Card 2:1 para a prévia social do GitHub e posts com link (exportar em 1280x640).
\documentclass{article}
\usepackage[paperwidth=25.6cm,paperheight=12.8cm,margin=0pt]{geometry}
% Macros comuns aos cards de divulgação — marca for_code.
\usepackage{fontspec}
\usepackage{xcolor}
\usepackage{graphicx}
\usepackage{tikz}
\usetikzlibrary{calc}
\newcommand{\fontpath}{../docs/assets/fonts/}
\setmainfont[Path=\fontpath, UprightFont=*-Regular, BoldFont=*-ExtraBold]{Lexend}
\newfontfamily\lexendlight[Path=\fontpath, UprightFont=*-Light]{Lexend}
\newfontfamily\lexendsemibold[Path=\fontpath, UprightFont=*-SemiBold]{Lexend}
\newfontfamily\lexendblack[Path=\fontpath, UprightFont=*-Black]{Lexend}
\setmonofont{DejaVu Sans Mono}[Scale=0.9]

\definecolor{darkpurple}{HTML}{1E1647}
\definecolor{lightpurple}{HTML}{614AD3}
\definecolor{yellow}{HTML}{FFFF00}
\definecolor{codebg}{HTML}{140E30}

\pagestyle{empty}
\hyphenpenalty=10000
\exhyphenpenalty=10000

\newcommand{\logo}[1]{\includegraphics[width=#1]{../docs/assets/logos/for_code_para_fundo_escuro.png}}
\newcommand{\simbolo}[1]{\includegraphics[width=#1]{../docs/assets/logos/simbolo_para_fundo_escuro.png}}
% Título com os colchetes da marca: \tit{texto}
\newcommand{\tit}[1]{\textcolor{lightpurple}{[}\,#1\,\textcolor{lightpurple}{]}}
% \numero{coordenada}{número}{rótulo}
\newcommand{\numero}[3]{%
  \node[anchor=north west, text=yellow, font=\lexendblack\fontsize{40}{40}\selectfont] at (#1) {#2};
  \node[anchor=north west, text=white, font=\lexendlight\small, text width=3.9cm, align=left] at ($(#1)+(0.05,-1.55)$) {#3};}

\newcommand{\kw}[1]{\textcolor{lightpurple}{\textbf{#1}}}
\begin{document}
\begin{tikzpicture}[remember picture, overlay]
  \fill[darkpurple] (current page.south west) rectangle (current page.north east);
  \node[anchor=center, opacity=0.15] at ([xshift=7.3cm]current page.center) {\simbolo{8.5cm}};
  \node[anchor=north west, inner sep=0pt] at ([xshift=1.3cm,yshift=-1.1cm]current page.north west) {\logo{3.8cm}};
  \node[anchor=north west, text=white, inner xsep=0pt, font=\bfseries\fontsize{46}{52}\selectfont] (t) at ([xshift=1.3cm,yshift=-3cm]current page.north west)
    {\tit{MySQL e SQL}};
  \node[anchor=north west, text=white, inner xsep=0pt, text width=11.8cm, align=left, font=\lexendlight\large] (s) at ([yshift=-3mm]t.south west)
    {Guia, slides e script da aula: do modelo relacional ao JOIN em duas horas, praticando direto no navegador.};
  \node[anchor=west, fill=codebg, draw=lightpurple, line width=1pt, rounded corners=3pt,
        inner sep=10pt, text=white, align=left, font=\ttfamily\normalsize] at ([xshift=15.2cm,yshift=0.2cm]current page.west)
    {\kw{SELECT} m.nome, c.nome\\
     \kw{FROM}\ \ \ membro m\\
     \kw{INNER JOIN} curso c\\
     \ \ \kw{ON} m.curso\_id = c.id;};
  \node[anchor=south west, text=yellow, font=\ttfamily\small] at ([xshift=1.3cm,yshift=0.9cm]current page.south west) {andradasdev.github.io/sql};
  \fill[yellow] (current page.south west) rectangle ([yshift=0.15cm]current page.south east);
\end{tikzpicture}
\end{document}
